\chapter{Tinjauan Pustaka}

\section{Penelitian Terdahulu}
\TemplateTodo{Ringkas penelitian terdahulu yang paling relevan. Jelaskan tujuan, metode, data, hasil, dan keterbatasannya. Akhiri dengan posisi penelitian Anda dibandingkan karya-karya tersebut.}

Tinjauan pustaka bukan sekadar daftar ringkasan artikel. Susun pembahasan secara tematik sehingga pembaca memahami perkembangan pengetahuan, kesenjangan yang masih ada, dan alasan penelitian ini diperlukan.

\begin{table}[H]
\centering
\caption{Contoh format ringkasan penelitian terdahulu}
\begin{tabularx}{\linewidth}{@{}p{0.18\linewidth}X X@{}}
\toprule
Peneliti & Fokus dan Metode & Keterkaitan dengan Skripsi \\
\midrule
Nama peneliti (tahun) & Ringkasan pendek metode, data, dan hasil utama. & Jelaskan pelajaran, celah, atau pembeda penelitian Anda. \\
\bottomrule
\end{tabularx}
\end{table}

\section{Landasan Teori}
\TemplateTodo{Jelaskan konsep, teori, algoritma, standar, atau kerangka kerja yang benar-benar digunakan pada penelitian.}

Setiap subbagian teori sebaiknya dihubungkan dengan kebutuhan metodologi. Hindari memasukkan teori yang tidak pernah digunakan pada Bab III atau Bab IV.

\subsection{Konsep Utama Pertama}
Tuliskan definisi, prinsip kerja, dan relevansi konsep terhadap penelitian. Gunakan sitasi untuk definisi atau teori yang bukan pendapat pribadi.

\subsection{Konsep Utama Kedua}
Tambahkan konsep lain sesuai kebutuhan. Jika menggunakan rumus, jelaskan arti setiap simbol dan alasan rumus tersebut digunakan.

\section{Kerangka Pemikiran}
\TemplateTodo{Jelaskan hubungan antara masalah, teori, metode, data, dan keluaran penelitian. Bagian ini dapat dibantu dengan gambar alur konseptual.}

Kerangka pemikiran membantu menunjukkan bahwa penelitian memiliki alur logis. Jika menambahkan gambar, simpan aset ke folder \texttt{figures/} dan gunakan perintah \texttt{\\includegraphics}.
